% !TeX root = ../papers/nc-p3.tex
\section{節点中心を持つ全特殊値のKummer類による分類}
\label{sec:nodal-center-special-values}
ここでは有限四辺の窓の、次の有限検査で定まる部分を扱う。
四つの中心四次パラメータについて全次数の正則性・自由性の有限切断検査を満たし、
全中心はその四元上の階数4の整な環で、その射影中心 $(X,H)$ は
$\mathbb P^4$ の四次超曲面、または滑らかな三次元二次曲面上の四次分岐二重被覆であるとする。
後者の中心環は重み $(1,1,1,1,1,2)$ の二次・四次完全交叉とする。
両方で特異点は通常二重点だけとする（初めにその数は課さない）。
全中心とその積はパラメータ環上の有限自由な経路表示で全ての矢との可換核を取って計算できる。
節点条件も射影的な特異集合とアフィンチャート上の二次主要部の階数で判定できる。

\begin{theorem}\label{thm:nodal-center-special-values}
この部分の各窓について $X$ は三十節点を持ち、そのうち二十点が薄い単純を持つ。
一般性、窓の属する成分、点配置の被約性・対称性、一般商の指数を仮定する必要はない。
$U=X_{\rm sm}$ とすれば
\[
 \operatorname{Pic}(U)=\mathbb Z H,\qquad
 \operatorname{Br}(U)=\mathbb Z/4,\qquad
 0\longrightarrow \langle H\bmod4\rangle\longrightarrow
 H^2_{\rm et}(U,\mu_4)\longrightarrow\operatorname{Br}(U)[4]\longrightarrow0 .
\]
添字保持同型類は、実現される偏極中心 $(X,H)$ と、右端において生成元に写る
中央群の元 $\beta$ との組の同型類で分類される。各実現中心に対して八つの
$\beta$ は全て実現される。同じ中心の異なる印の同一視は偏極自己同型の自然な作用だけである。
添字を一つ移す作用は符号規約を合わせて $\beta\mapsto\beta+H$ であり、
最小周期は $\beta+sH$ が $\beta$ の自己同型軌道に入る最小の $s=1,2,4$ である。
この記述は中心の型に属する全特殊値に適用できる。どの偏極中心が実現されるかの
無条件な判別と、節点以外を持つ中心は別の問題である。
\end{theorem}
\begin{proof}[一般性を使わない局所型]
$P$ を中心四パラメータの多項式環、$Z$ を全中心、$B$ を折り畳み環とする。
$P$ 上の自由性と $P_+$ で割ったFrobenius性から
$B\simeq\operatorname{Hom}_Z(B,Z)$ が左加群として成立する
（命題\ref{prop:smooth-azumaya}の証明の通り、両方の双対の次数差は一致する）。
一般分数体上の有限次元代数は従って自己入射的で、一方でKoszul両側分解から有限大域次元を持つ。
ゆえに半単純で、中心が体なので中心単純である。$P$ 上では十六個の隅が各階数64、
中心拡大は四次なのでPI次数は16、各頂点の一般行列階数は4である。
同命題の被約traceによる証明から滑らかな部分はAzumayaで、節点においてはAzumayaでない。
$H$ の切断を1にして次数方向を除くと、各単純の全頂点次元は同じ $d$ である。
これは定理\ref{thm:twenty-minimal-characters}のHilbert分子の四乗根での消滅による。
完備節点では既に証明したBrauer類の消滅により一般商は分裂する。
原始射影の一般行列階数を $m_j$ とすれば
$4=\sum d_jm_j$、かつ格子の有限射影分解から $\gcd_j(m_j)=1$ である。
この二条件はある $m_j=1$ を強制する（全て2以上なら4の分割で公約数が残る）。

その原始元 $p$ を取ると $pBp=R$ で、完備代数は高さ1からの延長により
$\operatorname{End}_R(M)$、$M=Bp$ である。$M$ は $R$ 直和因子を持つ極大Cohen--Macaulay加群である。
三次元通常二重点の非自由な直既約極大Cohen--Macaulay加群は二つの罫線の階数1加群 $I,I^*$ である。
ここではKn\"orrer周期性を $uv$ と $uv-st$ の行列分解に適用する。
一次元の $uv=0$ の極大Cohen--Macaulay加群の分類は、正規化上の二自由格子に写し、
導手で割った二ベクトル空間に両射影が全射である部分空間を取ることでも分かる。
各側だけの部分空間を除くと同型のグラフになり、直既約は自由加群と二枝の二加群だけである。
二枝の一行一列行列分解を周期移動すると二つの標準二行二列表示を得る。
$I,I^*$ は $M$ に同時には現れない。間のHomは局所類群の $\pm2$ の階数1加群で、
上の分類から極大Cohen--Macaulayでないが、$\operatorname{End}_R(M)$ の全直和因子は極大Cohen--Macaulayである。
非Azumaya性から $M$ は自由だけでもない。従って二原始型の階数は全て1、
単純の頂点次元は $1,3$ または $2,2$ である。

各型の単純には、次数を自明化した後に自己第一拡大がない。
基本代数の対角根基は $R$ の極大イデアルであり、これは
$I^*I=(u,v,s,t)$ すなわち非対角二因子の積に入るからである。
薄い四辺の表現は有限点グラフの四組と一対一で、中心の座標値を1にした変形においても
同じ対応がある。各組で非零の中心座標を取れば、四辺の倍率の積はその値により定まり、
残りは頂点基底変更で消える。従って自己拡大の消滅により全ての点は被約である。
長さ20の点グラフ上の二十薄点はそれぞれ異なる $1,3$ 型の節点を定める
（一節点には薄い単純は一種類だけである）。
\end{proof}

\begin{proof}[不分岐群と因子]
前節のねじれなしの局所Euler数の計算はこの全ての局所模型に適用でき、
滑らかな変形のEuler数は両中心型で $-56$ である。二重被覆の場合は
$2\chi(Q^3)-\chi(Q^3\cap F_4)=8-64$ である。
分岐曲面の第二Chern類は $(1+h)^5/((1+2h)(1+4h))$ の二次項 $8h^2$ である。
従って節点数は $4=-56+2N$ より30となる。
全節点の吹上げ $Y$ 上に一般のAzumaya類が延長し、定理\ref{thm:thirty-node-dihedral-cover}の
長さ64のねじれ充満例外列の証明が各点の上でそのまま働く。
その証明の $K^1$ の消滅により $H^3(Y,\mathbb Z)$ は位相的ねじれ類で生成される有限巡回群である。
また $H^1(Y,\mathcal O)=H^2(Y,\mathcal O)=0$ である。節点は有理特異点で
$X$ 上のこの消滅は超曲面の列、二重被覆の場合は
$\mathcal O_{Q^3}\oplus\mathcal O_{Q^3}(-2)$ から従う。
$Y$ のEuler数は $-56+30\chi(\mathbb P^1\times\mathbb P^1)=64$ であるから、
Poincar\'e双対と $b_1=b_3=0$ より $b_2=31$ となる。
$H$ と三十の例外因子は独立である。各例外の罫線に交差させればその例外の係数が決まり、
$H^3=4$ が残る係数を決める。よって $\operatorname{Pic}(U)\otimes\mathbb Q=\mathbb Q H$ で、
局所自由巡回類群への像は全て零である。従って $U$ の因子の閉包は各節点でCartierである
（反射的階数1層の完備化が自由なら元も自由）。$\operatorname{Pic}(U)=\operatorname{Pic}(X)$ を得る。
この群は捩れを持たない。実際両型の滑らかな変形は滑らかなFano三次元多様体なので
単連結である。節点のリンクもMilnor繊維も単連結であり、特異繊維と滑らかな繊維は
節点の小球の外で同相なのでvan Kampenにより $X$ は単連結である。
指数列と普遍係数定理から $\operatorname{Pic}(X)$ に捩れはない。また一般超平面切断のK3上で
$H$ の平方は4で格子は偶なので $H$ は原始である。これでPicard群が決まった。

薄い節点が存在するので定理\ref{thm:fixed-brauer-center-windows}の局所行列式の奇偶により
$\alpha$ の周期は4である。次数4の隅から指数も4となる。指数列と上の $H^3$ の巡回性から
$\operatorname{Br}(U)=\langle\alpha\rangle$ も従う。以上の複素数上の証明は
滑らかな射影解消上の有限係数比較と局所Picard群を介して幾何学的標数零で成立する。
\end{proof}

\begin{proof}[八つの印と特殊自己同型]
$U$ の分裂gerbe上の零番束について、その行列式を自明化した分裂の $\mu_4$-gerbeの類を取る。
双対規約と併せ、その符号は零番束の行列式に通常の束 $L$ を掛けたとき
Kummer類が $+[L]$ だけ変わる規約で $\beta$ とする。分裂束の全てを一つの通常の束で
捩っても行列式の変化は四倍なのでこの印は変わらない。
超平面上のHilbert多項式の計算により、隣り合う二束の行列式の比は $H$ である。
同じ生成元を持つ二窓について、定理\ref{thm:fixed-brauer-center-windows}の
一般超平面の定義体上の安定性・Riemann--Rochの証明が適用できる。
二重被覆でも普遍超平面は $U$ 上の射影空間束で、その一般繊維は滑らかなK3である。
従って $\beta$ が一致すれば、四移動を除いて四束が一般切断上で周期的に同型になり、
全矢と二次の関係の同型が得られる。逆も行列式の構成から従う。

ここで得た同型が指定した中心同定を保つことを、特殊中心の自己同型を含めて確かめる。
四束の周期的同型は全ての長さの積の切断上の評価を運び、スカラー切断はそのまま運ぶ。
従って中心自己同型は一般超平面切断を点ごとに固定する。$|H|$ に誘導される
射影変換はその超平面を点ごとに固定するため、行列として $1+v\ell$ の形である。
$\ell$ は一般超平面の方程式とする。この変換は標数零の定数点である三十の節点像を置換する。
二重被覆でも節点は分岐上にあり、それぞれの像は相異なる。この集合は一直線に入らない。
四次超曲面では三次の全偏微分がその直線で消えてしまう。
二重被覆では二次曲面が直線を含み、その上で枝の四次式と二式のJacobianの
四次小行列式（ここでは二行二列の次数4の小行列）が全て恒等的に消えてしまう。

もし変換が恒等でなければ、各移動点とその像の直線は $[v]$ を通る。
固定点は $[v]$ だけである（一般超平面は節点を通らない）。
非共線性から異なる二つの移動直線があるので、その交点 $[v]$ は定数点である。
定数ベクトル $v_0$ を選んで $v=cv_0$ とすると、移動する定数点 $p$ と像から
$c\ell(p)$ は定数に決まる。異なる二つの移動点でその比を取ると
一般の $\ell$ における二評価の比が定数となり矛盾する。よって射影変換は恒等である。
二重被覆の非自明な被覆変換は一般切断を点ごとに固定しないのでこれも除かれる。
関数体上の同型は有限型の同型方程式により代数閉定数体上の同型を与える。

Kummer列は各生成元に四つの印を与え、移動でその四つを尽くす。反対窓で
もう一つの生成元を得るため、実現中心上では全八点が実現される。
任意の偏極同型で印を引き戻して上の中心固定判定を適用すれば、同型条件と周期条件が従う。
\end{proof}
